\documentclass[10pt,a4paper]{amsart}
\usepackage[margin=19mm]{geometry}
\usepackage{amsmath,amssymb,mathtools}
\usepackage[scr=rsfs,cal=euler]{mathalfa}
\usepackage{xfrac}
\usepackage[unicode,hidelinks,bookmarks=false]{hyperref}
\newcommand{\ie}{i.e.\ }
\newcommand{\cf}{cf.\ }
\newcommand{\resp}{respectively\ }
\newcommand{\Subsection}{\subsection}
\providecommand{\nobreakdash}{\nobreak\hbox{-}}
\setlength{\emergencystretch}{2em}
\multlinegap0pt
\usepackage{bm}
\usepackage{bbm}
\usepackage{dsfont}
\usepackage{xspace}
\usepackage{cancel}
\usepackage{mathtools}
\usepackage{amsthm}
\makeatletter
\@ifundefined{assumption}{\newtheorem{assumption}{Assumption}}{}
\@ifundefined{lemma}{\newtheorem{lemma}{Lemma}}{}
\@ifundefined{theorem}{\newtheorem{theorem}{Theorem}}{}
\makeatother
\newcommand{\lspan}{\mathrm{span}}
\title[Perturbation theory for killed Markov processes and quasi-st]{Perturbation theory for killed Markov processes and quasi-stationary distributions}
\author{Daniel Rudolf, Andi Q. Wang}
\date{Source: arXiv:2109.13819v2 (CC BY 4.0)}
\begin{document}
\maketitle

\noindent\emph{Source and licence.} Perturbation theory for killed Markov processes and quasi-stationary distributions. Version arXiv:2109.13819v2 (CC BY 4.0), distributed under the Creative Commons Attribution 4.0 International licence, as recorded in the arXiv licence metadata for this exact version and shown on the abstract page https://arxiv.org/abs/2109.13819v2. Full rights notice: licenses/arxiv-2109.13819-CC-BY-4.0.txt.\par\smallskip

\noindent\textit{Editorial scope.} Counted unit: the Theorem whose source label is \texttt{eq\_thm\_bnd\_case} in the article (source0.tex lines 351--359, source label \texttt{eq\_thm\_bnd\_case}), delivered together with the article's own complete original proof of it (source0.tex lines 360--418, one \texttt{proof} environment of 59 lines). No mathematics is written, repaired or re-derived: statement and proof are the article's own text line for line. Required context retained uncounted: assm:spec\_gap (lines 269--272); lemma:wey (lines 303--310); assm:H\_small (lines 313--316). Every definition, display and label of those lines is retained unchanged. Omitted from this package and not redistributed: 1--268, 273--302, 311--312, 317--350, 419--1561. Nothing inside a retained excerpt is omitted or altered: every retained body is delivered exactly as the article's lines are written, so no omission receipt of any kind accompanies this package. Citations: the retained excerpts carry no citation command at all, so no bibliography entry of the article is transcribed and no reference is redistributed. Reference adaptation: none was needed and none was made; every reference inside the delivered unit resolves inside it, so no unresolved reference or citation is produced. Printed slips kept literal: no printed word, symbol, hypothesis or number of the counted statement or of its proof was altered. Layout and class: the article's own class is \texttt{aptpub} with packages including babel, inputenc, subcaption, commath, graphicx, color; the wrapper is the lane's pinned \texttt{amsart} editorial wrapper at 10pt on A4, which restates the theorem-like environments the retained text needs and adds the article's own macro definitions that the retained text needs (\texttt{\textbackslash lspan}), with extra wrapper packages (bm, bbm, dsfont, xspace, cancel, mathtools). The delivered numbering is the unit's own. Assets: no retained span contains a figure environment, a graphics-inclusion command, a PDF-page inclusion, a table or a TikZ picture; no image file is copied into this package, the package declares no asset dependency and no image file is redistributed with it. Licence: version arXiv:2109.13819v2 of this article is distributed under the Creative Commons Attribution 4.0 International licence, as recorded in the arXiv licence metadata for this exact version and shown on the abstract page https://arxiv.org/abs/2109.13819v2. A full rights notice is delivered at licenses/arxiv-2109.13819-CC-BY-4.0.txt. Characters: every retained excerpt is pure ASCII, so the delivered engine, XeLaTeX with its default Latin Modern text font, reports no missing character.
\begin{assumption}
	For the operator $L$ we have that $\nu := \lambda_1-\lambda_0 >0$, that is, $L$ possesses a spectral gap. Additionally, we assume that the eigenspace corresponding to $\lambda_0$ is $1$-dimensional.
    \label{assm:spec_gap}
\end{assumption}

\begin{lemma}[Weyl's lemma]
	\label{lemma:wey}
    For $j=0,1$, we have 
    \begin{equation}
        |\lambda_j - \widehat \lambda_j|\le \|H\|.
        \label{eq:eval_bd}
    \end{equation}
\end{lemma}

\begin{assumption}
    Suppose that the operator norm of $ H$ satisfies $\|H\|<\nu/2$, where $\nu = \lambda_1-\lambda_0$ is the spectral gap; see Assumption~\ref{assm:spec_gap}.
    \label{assm:H_small}
\end{assumption}

\begin{theorem}
    Under Assumption~\ref{assm:spec_gap} and~\ref{assm:H_small} we have
    \begin{equation} \label{eq_thm_bnd_case}
        \|\varphi-\widehat \varphi\| 
        % \le \frac{\| H\varphi-\langle H\varphi,\varphi \rangle \varphi \|}{\nu-2\|H\|}
        \leq \frac{\| H\varphi \|}{\nu-2\|H\|}.
    \end{equation}
    \label{thm:bdd_case}
\end{theorem}

\begin{proof}
Recall that $\rho := \langle \widehat \varphi, \varphi\rangle \ge 0$ and with $\mathbb{G}:=\lspan\{ \varphi\}$ define
	\begin{equation*}
	h := \frac{P_{\mathbb{G}^\perp} \widehat{\varphi}}{\Vert P_{\mathbb{G}^\perp} \widehat{\varphi} \Vert}
	= (1-\rho^2)^{-1/2} (\widehat \varphi - \rho\, \varphi).
	\end{equation*}
	(If $\rho=1$, then $\widehat \varphi = \varphi$ and the result trivially holds true.) Note that $\|h\|=1$
	and $h\perp \varphi$ as well as $L h \perp \varphi$ by the self-adjointness of $L$. 
 Consequently, by using the representation of $h$, $L\varphi = \lambda_0 \varphi$ and $\widehat L=L+H$ we have that
 \begin{align*}
  \rho \lambda_0 \varphi + \rho H \varphi + \sqrt{1-\rho^2} \widehat L h =
  \rho \widehat L \varphi + \widehat L (\widehat \varphi -\rho \varphi)
  = 
    \widehat L\widehat\varphi
= \widehat{\lambda}_0 \widehat\varphi.
 \end{align*}
Taking the inner product with $h$, employing $\langle \widehat \varphi,h \rangle = \sqrt{1-\rho^2}$, $\varphi\perp h$ and rearranging yields
\begin{equation}
\label{eq: repres}
    -\rho \langle H \varphi, h\rangle = \sqrt{1-\rho^2} \left( \langle h,\widehat L h\rangle -\widehat\lambda_0 \right).
\end{equation}
Moreover, by $\widehat{L}=L+H$, $h\perp \varphi$ and Lemma~\ref{lemma:wey} we have 
\begin{align*}
   B & := \langle h,\widehat{L} h \rangle - \widehat\lambda_0
    = \langle h,L h \rangle + \langle h,H h \rangle - \widehat\lambda_0 
    \geq \lambda_1 + \langle h,H h \rangle - \widehat\lambda_0  \\
    &\geq \lambda_1 - \Vert H \Vert - \widehat\lambda_0
    \geq \lambda_1 - 2\|H\|-\lambda_0=\nu-2\|H\|,
\end{align*}
which is, due to Assumption~\ref{assm:H_small}, strictly positive. 
Squaring \eqref{eq: repres}, rearranging, and taking a square root, with $A:=|\langle H \varphi ,h\rangle|$ we obtain
\begin{equation*}
    \rho = \frac{B}{\sqrt{A^2 + B^2}}.
\end{equation*}
Then
\begin{equation*}
    1-\rho = \frac{\sqrt{A^2+B^2}-B}{\sqrt{A^2+B^2}} = \frac{A^2}{(\sqrt{A^2+B^2}+B)\sqrt{A^2+B^2}}\le\frac{A^2}{2B^2},
\end{equation*}
from which we have that $\sqrt{2(1-\rho)}\le A/B$.
The fact that we can represent $\widehat \varphi$ by
	the orthogonal decomposition
	\begin{equation*}
	\widehat \varphi
	= P_{\mathbb{G}} \widehat \varphi +P_{\mathbb{G}^\perp} \widehat \varphi
	=
	\rho\, \varphi + \sqrt{1-\rho^2 } h
	\end{equation*}
 gives, with $\rho = \langle \widehat \varphi, \varphi\rangle \ge 0$, that
$
\widehat \varphi - \varphi = (\rho-1)\varphi + \sqrt{1-\rho^2} h.
$
Consequently, 
\[
\Vert \widehat\varphi -\varphi \Vert = \sqrt{2(1-\rho)} \leq \frac{A}{B} \leq \frac{\Vert H\varphi \Vert}{\nu - 2\Vert H \Vert},
\]
which finishes the proof.


\end{proof}

\end{document}
